% -*- Document-Type: LaTeX; Master-File: manual.tex -*- 

%\chapter{Contributed works}
% [ck: Aug 16 95]

\section{Description of the \elan\ contributions}


Many computational processes in Automated Deduction and Constraint
Programming can be expressed as instances of a general approach that
consists of applying transformation rules on formulas with some strategy, 
until reaching specific normal forms. Such processes
are naturally modelled in \elan, and can be classified according to
the area of interest, namely programming, proving and solving.

\noindent
{\it Programming}:
One of the first application was to prototype the fundamental
mechanisms of logic and functional programming languages  like first-order
resolution and $\lambda$-calculus.
%
The general framework of Constraint Logic
Programming can be easily designed 
in the \elan\ framework~\cite{KirchnerR-FI98}, since its operational
semantics is clearly formalised as rewrite rules, although the
application strategy is often defined in an informal way. 
%
Some implementations~\cite{Borovansky-95} related to a 
calculus of explicit substitutions (the first-order rewrite system
$\lambda\sigma$ that mimics $\lambda$-calculus) open the way of
implementing higher-order logic programming languages via a
first-order setting. Another calculus of 
explicit substitutions based on the $\pi$-calculus is used to provide
a formal specification of Input/Output for \elan~\cite{Viry-RwLg1996}. 

\noindent
{\it Proving}:
\ELAN\ was used in order to implement a predicate prover based on the
rules proposed by J.-R.~Abrial, and implemented in the
B-tools~\cite{CirsteaKirchner97web}.
We developed also a propositional sequent calculus, completion procedures for
rewrite systems~\cite{KirchnerM-RTA95}, sufficient conditions
for the termination problem~\cite{genetgnaedig-caap97}.
A library for automata construction and manipulation has been
designed. Approximation automata are used to check conditions for
reachability, sufficient completeness, absence of conflicts in systems 
described by non-conditional rewrite rules~\cite{genet-rta98,genet-these}.

\noindent {\it Checking}: A very special case of proof is just exhaustive
check of all the possibilities. Using in particular ``dont know'' strategies, it
was efficiently rediscovered by exhaustive search that the well known Needam and
Schroeder authentification protocol is unsafe~\cite{Cirstea-Proto-99}.

%\subsubsection{Solving}
\noindent
{\it Solving}:
The notion of rewriting controlled by strategies is used 
in~\cite{Castro-FI98,KirchnerR-FI98} to describe in a unified way the
constraint solving 
mechanism as well as the meta-language needed to manipulate the
constraints~\cite{castro-esslli-1997,castrokirchner:rta-war:1998}. This provides programs that are
very close to the proof theoretical setting used now to describe
constraint manipulations like unification or numerical constraint
solving. \elan\ offers a constraint programming environment where
the formal description of a constraint solver is directly
executable. 
\elan\ has been tested 
on several examples of constraint solvers for various computation domains and
combinations like abstract
domains~\cite{KirchnerR-FI98,Ringeissen-RTA97} (term  
algebras) and more concrete ones (finite domains, integers).
In~\cite{Castro-ICTAI96,Castro-RwLg1996,Castro-ESSLLI97,castro:phd:1998},  
it is shown how to use computational systems as a general framework
for handling Constraint Satisfaction Problems (CSP for short). The
approach leads to the design in \elan\ of \textsf{COLETTE}, a solver
for constraints over integers and finite domains~\cite{castro:aisc:1998}. 
We also combine rules, constraints and strategies in order to deal with problems like
planning and scheduling~\cite{DuboisKirchner-Ercim99}.


\section{\elan's annotated bibliography}

\begin{description}

\item[
\cite{Elan-RwLg1996,Elan-ASF+SDF97,BorovanskyKKMR-WRLA98}
] 
General presentations of the different \elan\ system releases.


\item[
\cite{Strategies-RwLg1996,BorovanK-ASF+SDF97,Borovansky-WRLA98}
]
A description of the operational as well as denotational semantics of
the new \elan\ strategies, available since version 2.0. This allows
the user to define its own strategies as a computational system.


\item[
\cite{BKK-Fuji-98}
]
Towards a functional operational semantics of \elan.

\item[
\cite{BorovanskyJMR-WRLA98}
]
Some useful applications related to the \REF\ format.


\item[
\cite{BorovanskyCastro-WRLA98}
]
The facilities provided in \elan\ for the cooperation of constraint
solvers.

\item[
\cite{BoroThese}
]
A full description of the \elan\ strategy language. In french.

\item[
\cite{Borovansky-95}
] 
An \elan\ implementation of higher-order unification by
using a calculus of explicit substitutions.


\item[
\cite{Castro-RwLg1996,Castro-ICTAI96,Castro-ESSLLI97}
] 
The use of \elan\ to specify various constraint manipulation
algorithms for CSP's.

\item[
\cite{Castro-FI98,KirchnerR-FI98,Castro-ESSLLI97,castro:aisc:1998,castro:phd:1998}
]
The rule-based approach to implement Constraint Programming and
Constraint Solving Techniques. 

\item[
\cite{CirsteaKirchner-LivreFroCoS99,CirsteaKirchnerINRIA99}
]
The presentation of the rewriting calculus that provides a full
semantics to \elan.

\item[
\cite{CirsteaKirchner97web}
]
The specification of Abrial's B predicate prover.

\item[
\cite{dubois98b}
]
The strategy language is used to build plans.

\item[
\cite{DuboisKirchner-Ercim99}
]
A combination of rewriting rules, constraints and strategy language for planning problems.

\item[
\cite{KirchnerKV-MIT95}
] 
The first presentation of the general
ideas developed in \elan, with the definition of computational
systems (including the definition of strategies) and the application
of \elan\ to design and to prototype constraint programming languages.


\item[
\cite{Reflection-RwLg1996}
] 
An illustration of the reflective power of \elan\ and rewriting logic.

\item[
\cite{KirchnerM-RTA95}
] 
The \elan\ implementation of some rule-based completion
procedures.

\item[
\cite{Kirchner-JFPLC99}
]
An \ELAN\ tutorial.


\item[
\cite{MoreauK-ASF+SDF97,MoreauK-PLILP+ALP98,KirchnerM-98,MoreauThese99}
]
A description of the new compilation techniques developed for \elan, in
particular for associativity and commutativity


\item[
\cite{Ringeissen-RTA97}
]
The use of \elan\ to interface various unification algorithms
and to build combination algorithms for unification.

\item[
\cite{Viry-RwLg1996}
] 
An implementation of input/output in \elan\ via an explicit
substitution calculus for the $\pi$-calculus.

\item[ \cite{Vittek-RTA96} 
] 
A description of the first \elan\ compiler.

\item[
\cite{VittekThese}
] 
The first main version of the \elan\ system in full details, from
design to implementation. In French. 


\end{description}

